\documentclass[12pt]{article}

% --- Packages ---
\usepackage[utf8]{inputenc}
\usepackage[T1]{fontenc}
\usepackage[french]{babel}
\usepackage{amsmath, amssymb}
\usepackage{geometry}
% \usepackage[draft]{graphicx}
\usepackage{tikz}
\usepackage[most]{tcolorbox}
\usepackage{multicol}
\usepackage{lmodern}
\usepackage{pgfplots}
\pgfplotsset{compat=1.18}
\usepackage{mdframed}
\usepackage{array}
\usepackage[thinlines]{easytable}
\usepackage{tkz-tab}
\usepackage{bm}
\usepackage{xcolor}
\usepackage{sectsty}
\usepackage[utf8]{inputenc}
\usepackage[T1]{fontenc}
\usepackage[inline]{enumitem}
\usepackage[dvipsnames]{xcolor}
\usepackage{multido}
% --- Page ---
\geometry{margin=2cm}

% Footer
\usepackage{fancyhdr}
\pagestyle{fancy}
\cfoot{\thepage}
\rfoot{\textit{G.M.}}
\renewcommand{\headrulewidth}{0pt}
\renewcommand{\footrulewidth}{0pt}

% Pour les normes : 
\newcommand*{\norme}[1]{\left\lVert\overrightarrow{#1}\right\rVert}

\usetikzlibrary{shapes.misc}

\tikzset{cross/.style={cross out, draw=black, minimum size=2*(#1-\pgflinewidth), inner sep=0pt, outer sep=0pt},
%default radius will be 1pt. 
cross/.default={4pt}}

% mettre dans le préambule, après \usepackage{xcolor}
\usepackage{sectsty} % gère facilement la couleur des titres
\setlength{\parindent}{0pt}

\sectionfont{\color{red}\Large\bfseries}     % pour \section*
\subsectionfont{\color{red}\large\bfseries\hspace{1cm}}

\begin{document}

% ============================
%   Probabilités
% ============================

\newpage

\begin{center}
    {\LARGE \textbf{Arithmétique}}\\[4mm]
    {\large -- 2\up{nde} --}
\end{center}

\vspace{10pt}
\hrule
\vspace{0pt}

% ============================
%   Partie 1 : Rappels sur les probabilités
% ============================

\section*{\textcolor{red}{I/ Multiples et diviseurs}}

\vspace{5pt}

\begin{mdframed}[
    linewidth=1pt,
    roundcorner=5pt,
    linecolor=ForestGreen,
    innertopmargin=10pt,
    innerbottommargin=10pt
]
{\color{ForestGreen}
\textbf{\underline{Définition :}}}
Soit $a$ et $b$ deux entiers naturels.

\vspace{5pt}
On dit que $a$ est un multiple de $b$ s’il existe un entier $k$ tel que $\bm{a = kb}$.

\end{mdframed}

\vspace{5pt}
\underline{Remarque : } On dit que $b$ est un \_\_\_\_\_\_\_\_\_\_\_\_ de $a$.

\vspace{5pt}
\begin{mdframed}[
    linewidth=1pt,
    roundcorner=5pt,
    linecolor=RoyalBlue,
    innertopmargin=10pt,
    innerbottommargin=10pt
]
{\color{RoyalBlue}
\textbf{\underline{Exercice d'application n°1:} Montrer qu'un nombre est un multiple ou un diviseur.}}

\vspace{5pt}
Les affirmations suivantes sont-elles vraies ou fausses : 

\vspace{5pt}
\hspace{1cm}a) 36 est un multiple de 12 \hspace{2cm}b) 28 est un multiple de 8

\vspace{5pt}
\hspace{1cm}c) 6 est un diviseur de 54 \hspace{2.3cm}d) 7 est un diviseur de 24

\vspace{0.5cm}
{\color{black!45}\dotfill}

\vspace{0.5cm}
{\color{black!45}\dotfill}

\vspace{0.5cm}
{\color{black!45}\dotfill}

\vspace{0.5cm}
{\color{black!45}\dotfill}

\vspace{20pt}

\end{mdframed}

\begin{mdframed}[
    linewidth=1pt,
    roundcorner=5pt,
    linecolor=Red,
    innertopmargin=10pt,
    innerbottommargin=10pt
]
{\color{red}
\textbf{\underline{Propriété :}}
La somme de deux multiples d’un entier $a$ est un multiple de $a$.
}
\end{mdframed}

\begin{mdframed}[
    linewidth=1pt,
    roundcorner=5pt,
    linecolor=Plum,
    innertopmargin=10pt,
    innerbottommargin=10pt
]
{\color{Plum}
\textbf{\underline{Démonstration :}}}

\vspace{5pt}
Soit $n \in \mathbb{N}$.

\vspace{5pt}
Soient $b$ et $c$ deux multiples de $n$.
Démontrons que la somme de deux multiples de $n$ est un multiple de $n$.

\vspace{5pt}
Comme $b$ est un multiple de $n$, il existe un entier $k_1$ tel que $b = 3k_1$.
\\Comme $c$ est un multiple de $n$, il existe un entier $k_1$ tel que $c = 3k_2$.

\vspace{5pt}
Alors : $b+c =$

\vspace{0.5cm}
{\color{black!45}\dotfill}

\end{mdframed}

\begin{mdframed}[
    linewidth=1pt,
    roundcorner=5pt,
    linecolor=RoyalBlue,
    innertopmargin=10pt,
    innerbottommargin=10pt
]
{\color{RoyalBlue}
\textbf{\underline{Exercice d'application n°2:} Résoudre un problème avec des multiples et des diviseurs.}}

\vspace{5pt}
Montrer que la somme de trois entiers consécutifs est toujours un multiple de 3.

\vspace{0.5cm}
{\color{black!45}\dotfill}

\vspace{0.5cm}
{\color{black!45}\dotfill}

\vspace{0.5cm}
{\color{black!45}\dotfill}

\vspace{20pt}

\end{mdframed}


\vspace{10pt}
\section*{\textcolor{red}{II/ Nombres pairs, nombres impairs}}

\begin{mdframed}[
    linewidth=1pt,
    roundcorner=5pt,
    linecolor=ForestGreen,
    innertopmargin=10pt,
    innerbottommargin=10pt
]
{\color{ForestGreen}
\textbf{\underline{Définition :}}}

\vspace{5pt}
Un nombre \textbf{pair} est un multiple de 2.

\vspace{5pt}
Un nombre \textbf{impair} est un nombre qui n'est pas pair.

\end{mdframed}

\vspace{5pt}
\textbf{On en déduit la propriété suivante :}

\begin{mdframed}[
    linewidth=1pt,
    roundcorner=5pt,
    linecolor=Red,
    innertopmargin=10pt,
    innerbottommargin=10pt
]
{\color{red}
\textbf{\underline{Propriété :}}

\vspace{0.5cm}
{\color{black!45}\dotfill}

\vspace{0.5cm}
{\color{black!45}\dotfill}

\vspace{20pt}

}
\end{mdframed}

\underline{Exemples :}

\vspace{5pt}
\textbullet \ $34 = 2k$ avec $k =$

\vspace{5pt}
\textbullet \ $57 = 2k + 1$ avec $k = $

\vspace{10pt}
\begin{mdframed}[
    linewidth=1pt,
    roundcorner=5pt,
    linecolor=Red,
    innertopmargin=10pt,
    innerbottommargin=10pt
]
{\color{red}
\textbf{\underline{Propriété :}}

\vspace{5pt}
De manière abrégée (que l'on peut retrouver facilement par le calcul) :

\vspace{-20pt}
\begin{center}
    \vspace{20pt}
    \textbf{pair + pair $\longrightarrow$ \_\_\_\_\_\_\_\_}

    \vspace{20pt}
    \textbf{pair + pair $\longrightarrow$ \_\_\_\_\_\_\_\_}

    \vspace{20pt}
    \textbf{pair + pair $\longrightarrow$ \_\_\_\_\_\_\_\_}

    \vspace{20pt}
    \textbf{pair $\times$ \_\_\_\_\_\_\_\_ $\longrightarrow$ pair}

    \vspace{20pt}
    \textbf{pair $\times$ \_\_\_\_\_\_\_\_ $\longrightarrow$ impair}
\end{center}

}
\end{mdframed}

\begin{mdframed}[
    linewidth=1pt,
    roundcorner=5pt,
    linecolor=Red,
    innertopmargin=10pt,
    innerbottommargin=10pt
]
{\color{red}
\textbf{\underline{Propriété :}} Le carré d'un nombre impair est impair.
}
\end{mdframed}

\begin{mdframed}[
    linewidth=1pt,
    roundcorner=5pt,
    linecolor=Plum,
    innertopmargin=10pt,
    innerbottommargin=10pt
]
{\color{Plum}
\textbf{\underline{Démonstration :}}}

\vspace{5pt}
\vspace{0.5cm}
{\color{black!45}\dotfill}

\vspace{0.5cm}
{\color{black!45}\dotfill}

\vspace{0.5cm}
{\color{black!45}\dotfill}

\vspace{20pt}

\end{mdframed}

\begin{mdframed}[
    linewidth=1pt,
    roundcorner=5pt,
    linecolor=RoyalBlue,
    innertopmargin=10pt,
    innerbottommargin=10pt
]
{\color{RoyalBlue}
\textbf{\underline{Exercice d'application n°3:} Un problème avec des nombres pairs ou impairs.}}

\vspace{5pt}
Montrer que le produit de deux entiers consécutifs est un nombre pair.

\vspace{0.5cm}
{\color{black!45}\dotfill}

\vspace{0.5cm}
{\color{black!45}\dotfill}

\vspace{0.5cm}
{\color{black!45}\dotfill}

\vspace{0.5cm}
{\color{black!45}\dotfill}

\vspace{0.5cm}
{\color{black!45}\dotfill}

\vspace{0.5cm}
{\color{black!45}\dotfill}

\vspace{20pt}

\end{mdframed}

\vspace{10pt}
\section*{\textcolor{red}{III/ Nombres premiers}}

\begin{mdframed}[
    linewidth=1pt,
    roundcorner=5pt,
    linecolor=ForestGreen,
    innertopmargin=10pt,
    innerbottommargin=10pt
]
{\color{ForestGreen}
\textbf{\underline{Définition :}}}

\vspace{5pt}
Un nombre est \textbf{premier} s'il possède exactement {\color{black!45}\dotfill}

\end{mdframed}

\vspace{20pt}
\underline{Remarques :}

\begin{itemize}
    \item Il existe un nombre \textbf{infini} de nombre premier.
    Mais nous n'allons pas le démontrer ici, ça serait un peu compliqué. Mais vu que vous êtes forts, ok on va le 
    faire.

    \begin{mdframed}[
        linewidth=1pt,
        roundcorner=5pt,
        linecolor=Plum,
        innertopmargin=10pt,
        innerbottommargin=10pt
    ]
    {\color{Plum}
    \textbf{\underline{Démonstration :}}}
    NE ME REMERCIEZ PAS JE SAIS QUE ÇA VOUS FAIT PLAISIR

    \vspace{5pt}
    Nous allons démontrer cela \textbf{par l'absurde} : partir d'une hypothèse, aboutir à une absurdité, donc forcément 
    notre hypothèse de départ est \textbf{fausse}.

    \vspace{5pt}
    Supposons donc qu'il existe un nombre \textbf{fini} de nombres premiers.
    \\Notons-les par exemple 
    : $p_1, p_2, ... 
    p_n$, ou $p_n$ est le plus grand nombre premier.

    \vspace{5pt}
    Construisons un nouveau nombre comme le produit de ces nombres premiers auquel on ajoute 1 : 

    \vspace{-20pt}
    $$
    N = p_1 \times p_2 \times ... \times p_n + 1
    $$

    \textbullet \ $N > 1$ ce qui implique qu'il admet au moins un diviseur premire (théorème fondamental de 
    l'arithmétique), comme vous l'avez vu dans les décompositions en facteurs premiers.

    \textbullet \ Or, $N$ n'est divisible par \textbf{aucun} des $p_i$ (car le reste de la division serait 1 !) :

    \vspace{0.5cm}
    {\color{black!45}\dotfill}
    
    \vspace{0.5cm}
    {\color{black!45}\dotfill}
    
    \vspace{20pt}

    \textbullet \ Donc $N$ admet un diviseur premier \textbf{différent} de tous les $p_i$.

    \vspace{5pt}
    Ce qui est \textbf{absurde}.
    
    \end{mdframed}

    \vspace{10pt}
    \item Le nombre 1 n'est pas premier car {\color{black!45}\dotfill}
    
\end{itemize}

\vspace{10pt}
\begin{mdframed}[
    linewidth=1pt,
    roundcorner=5pt,
    linecolor=Red,
    innertopmargin=10pt,
    innerbottommargin=10pt
]
{\color{red}
\textbf{\underline{Théorème :}}

\vspace{5pt}
Tout nombre non premier peut se décomposer en produit de facteurs premiers.
\\L’ordre des facteurs n’a pas d’importance.
}
\end{mdframed}

\begin{mdframed}[
    linewidth=1pt,
    roundcorner=5pt,
    linecolor=ForestGreen,
    innertopmargin=10pt,
    innerbottommargin=10pt
]
{\color{ForestGreen}
\textbf{\underline{Définition :}}}

\vspace{5pt}
 On dit qu’une fraction est irréductible, lorsque son numérateur et son
dénominateur n’ont pas de diviseur commun autre que 1.

\end{mdframed}




\end{document}
